\documentclass[11pt,a4paper]{article}
\usepackage[utf8]{inputenc}
\usepackage[vietnamese]{babel}
\usepackage{amsmath,amssymb}
\usepackage{geometry}
\geometry{a4paper, total={180mm,260mm}, left=15mm, right=15mm, top=18mm, bottom=18mm}

\usepackage{tikz}
\usetikzlibrary{shapes.geometric, arrows.meta, positioning, calc}
\tikzset{
  arr/.style={-{Stealth[scale=0.8]}},
  every node/.append style={align=center}
}
\usepackage{tcolorbox}
\tcbuselibrary{skins,breakable}
\usepackage{fontawesome}
\usepackage{fancyhdr}
\usepackage{enumitem}

\pagestyle{fancy}
\fancyhf{}
\renewcommand{\headrulewidth}{0.5pt}
\renewcommand{\footrulewidth}{0.5pt}
\fancyhead[L]{\small\bfseries\color{violet!85!black}\faBook\ SỔ TAY KIẾN THỨC TOÁN 2}
\fancyhead[R]{\small\bfseries\color{violet!85!black}Lớp Toán Thầy Lĩnh -- loptoanthaylinh.vn}
\fancyfoot[C]{\small\color{gray}-- Sổ tay kiến thức: Bài 36. Ôn tập chung học kì 1 --}

\setlength{\parindent}{0pt}
\setlength{\parskip}{0pt}

\begin{document}

% --- BANNER BÀI HỌC ---
\begin{tcolorbox}[
  enhanced,
  colback=violet!8!white,
  colframe=violet!70!black,
  arc=3mm, boxrule=1.2pt,
  left=3mm, right=3mm, top=2mm, bottom=2mm,
  halign=center
]
  {\LARGE\bfseries\color{violet!90!black}\faBookmark\ BÀI 36: ÔN TẬP CHUNG HỌC KÌ 1}\\[1mm]
  {\itshape\footnotesize\color{black!75} Đại tổng kết toàn diện 4 trụ cột kiến thức Toán 2 Học kì 1: Số học, Phép tính, Hình học và Đo lường trên 1 trang A4}
\end{tcolorbox}

\vspace{-2mm}

\noindent
\begin{minipage}[t]{0.49\textwidth}\vspace{0pt}
\begin{tcolorbox}[
  enhanced,
  colback=violet!4!white,
  colframe=violet!75!black,
  arc=2mm, boxrule=0.8pt,
  left=2mm, right=2mm, top=1.8mm, bottom=1.8mm,
  title={\footnotesize\bfseries\color{white}\faKey\ 1. Trụ cột Số học \& Phép tính}
]
{\footnotesize
\begin{itemize}\setlength{\itemsep}{1pt}\setlength{\leftskip}{-2.5mm}
  \item \textbf{Số đến 100:} Đọc, viết, cấu tạo số $a0 + b$, so sánh số chục trước rồi đến đơn vị.
  \item \textbf{Cộng trừ qua 10 phạm vi 20:} Thuộc các bảng cộng $9+, 8+, 7+, 6+$ và bảng trừ tương ứng.
  \item \textbf{Cộng trừ có nhớ phạm vi 100:}
  \begin{itemize}\setlength{\itemsep}{0.5pt}\setlength{\leftskip}{-2.5mm}
    \item Cộng: Đơn vị nhớ sang chục (thêm 1).
    \item Trừ: Đơn vị mượn 1 chục (bớt 1 ở chục).
  \end{itemize}
  \item \textbf{Tìm thành phần:} $x = \text{Tổng} - \text{Số hạng}$; $x = \text{Hiệu} + \text{Số trừ}$; $x = \text{Số bị trừ} - \text{Hiệu}$.
\end{itemize}
}
\end{tcolorbox}
\end{minipage}
\hfill
\begin{minipage}[t]{0.49\textwidth}\vspace{0pt}
\begin{tcolorbox}[
  enhanced,
  colback=teal!4!white,
  colframe=teal!75!black,
  arc=2mm, boxrule=0.8pt,
  left=2mm, right=2mm, top=1.8mm, bottom=1.8mm,
  title={\footnotesize\bfseries\color{white}\faSliders\ 2. Trụ cột Hình học, Đo lường \& Giải toán}
]
{\footnotesize
\begin{itemize}\setlength{\itemsep}{1pt}\setlength{\leftskip}{-2.5mm}
  \item \textbf{Hình học phẳng:}\\
  Điểm, đoạn thẳng, đường thẳng, đường cong, 3 điểm thẳng hàng; độ dài đường gấp khúc $=$ tổng các đoạn thẳng.
  \item \textbf{Đo lường:}\\
  $1\text{ dm} = 10\text{ cm}$; Ki-lô-gam (kg); Lít ($l$);\\
  $1\text{ ngày} = 24\text{ giờ}$; $1\text{ giờ} = 60\text{ phút}$.
  \item \textbf{Giải toán có lời văn:}\\
  Nhiều hơn $\to (+)$; Ít hơn $\to (-)$;\\
  Hơn kém bao nhiêu $\to (\text{Lớn} - \text{Bé})$.
\end{itemize}
}
\end{tcolorbox}
\end{minipage}

\vspace{-1.5mm}
\begin{tcolorbox}[
  enhanced,
  colback=orange!4!white,
  colframe=orange!80!black,
  arc=2mm, boxrule=0.8pt,
  left=2mm, right=2mm, top=1.5mm, bottom=1.5mm,
  title={\footnotesize\bfseries\color{white}\faSitemap\ 3. Sơ đồ tư duy: Cây đại thụ 4 nhánh kiến thức Toán 2 Học kì 1}
]
\centering
\begin{tikzpicture}[
  node distance=0.8cm,
  box/.style={rectangle, rounded corners=2pt, draw=violet!80!black, fill=violet!20, text width=2.4cm, align=center, font=\scriptsize\bfseries, inner sep=3pt},
  branch/.style={rectangle, rounded corners=2pt, draw=#1!80!black, fill=#1!10, text width=2.5cm, align=center, font=\scriptsize, inner sep=2.5pt},
  arr/.style={-{Stealth[scale=0.8]}, thick, #1!80!black}
]
  \node[box] (root) {ĐẠI THỤ KIẾN THỨC\\TOÁN 2 HK1};
  \node[branch=teal, above left=0.3cm and 0.4cm of root] (b1) {\textbf{SỐ HỌC \& PHÉP TÍNH}\\Số đến 100\\Cộng trừ có nhớ $\le 100$};
  \node[branch=orange, above right=0.3cm and 0.4cm of root] (b2) {\textbf{HÌNH HỌC PHẲNG}\\Điểm, đoạn, đường gấp khúc\\Đếm hình tam giác, tứ giác};
  \node[branch=green, below left=0.3cm and 0.4cm of root] (b3) {\textbf{ĐO LƯỜNG THỰC TẾ}\\cm, dm, kg, $l$\\Ngày, giờ, phút, tháng};
  \node[branch=purple, below right=0.3cm and 0.4cm of root] (b4) {\textbf{GIẢI TOÁN CÓ LỜI VĂN}\\Thêm/bớt, Nhiều/ít hơn\\Hơn kém nhau bao nhiêu};

  \draw[arr=teal] (root) -- (b1);
  \draw[arr=orange] (root) -- (b2);
  \draw[arr=green] (root) -- (b3);
  \draw[arr=purple] (root) -- (b4);
\end{tikzpicture}
\end{tcolorbox}

\vspace{-1.5mm}
\noindent
\begin{minipage}[t]{0.49\textwidth}\vspace{0pt}
\begin{tcolorbox}[
  enhanced,
  colback=green!4!white,
  colframe=green!60!black,
  arc=2mm, boxrule=0.8pt,
  left=2mm, right=2mm, top=1.8mm, bottom=1.8mm,
  title={\footnotesize\bfseries\color{white}\faCheckCircle\ 4. Ví dụ thực hành chuẩn mực}
]
{\footnotesize
\textbf{Bài toán tổng hợp HK1:}\\
Lớp 2A có $18$ bạn nữ và số bạn nam nhiều hơn số bạn nữ là $4$ bạn. Hỏi:\\
a) Lớp 2A có bao nhiêu bạn nam?\\
b) Lớp 2A có tất cả bao nhiêu học sinh?\\
\textit{Bài giải:}\\
a) Số bạn nam là: $18 + 4 = \mathbf{22}$ (bạn).\\
b) Cả lớp có tất cả là: $18 + 22 = \mathbf{40}$ (học sinh).\\
Đáp số: a) $22$ bạn nam; b) $40$ học sinh.
}
\end{tcolorbox}
\end{minipage}
\hfill
\begin{minipage}[t]{0.49\textwidth}\vspace{0pt}
\begin{tcolorbox}[
  enhanced,
  colback=red!4!white,
  colframe=red!75!black,
  arc=2mm, boxrule=0.8pt,
  left=2mm, right=2mm, top=1.8mm, bottom=1.8mm,
  title={\footnotesize\bfseries\color{white}\faExclamationTriangle\ 5. Mẹo tránh bẫy trắc nghiệm}
]
{\footnotesize
\begin{itemize}\setlength{\itemsep}{1pt}\setlength{\leftskip}{-2.5mm}
  \item \textbf{Bí quyết đạt điểm 10 thi HK1:}\\
  1. Đọc kĩ đề, gạch chân số và câu hỏi.\\
  2. Đặt tính thẳng hàng cẩn thận ra nháp.\\
  3. Kiểm tra số nhớ của phép cộng và phép trừ.\\
  4. Đổi về cùng đơn vị đo trước khi tính.\\
  5. Luôn ghi đáp số kèm tên đơn vị trong ngoặc.
\end{itemize}
}
\end{tcolorbox}
\end{minipage}


\end{document}
